

\section{Le traumatisme de Sedan : Le basculement du modèle Rhétorique (Jésuite/Production) au modèle Philologique (Allemand/Analyse)}

L'histoire de l'enseignement en France, particulièrement celle des humanités classiques, ne saurait
se lire comme une simple succession de réformes administratives ou d'ajustements pédagogiques. Elle
est, dans sa chair même, une histoire politique, intimement liée aux soubresauts de l'identité
nationale et aux traumatismes collectifs. Parmi ces ruptures, aucune n'est plus profonde, plus
viscérale, que la défaite de 1870 face à la Prusse. Le désastre de Sedan ne marqua pas seulement la
chute du Second Empire et la perte de l'Alsace-Lorraine ; il provoqua une onde de choc
intellectuelle que Claude Digeon a magistralement qualifiée de « crise allemande de la pensée
française ».\cite{I-1-1-ref1}

Au lendemain de la débâcle, une sentence s'imposa avec la force d'une évidence douloureuse au sein
de l'élite républicaine naissante : « C'est l'instituteur prussien qui a gagné la bataille de Sadowa
», et par extension, celle de Sedan. La France n'avait pas seulement été vaincue par le nombre ou la
stratégie militaire, mais par une supériorité intellectuelle et organisationnelle, celle de la
\textit{Wissenschaft} allemande. Dès lors, l'Université française, et en son cœur l'enseignement des
langues anciennes, se trouva placée au banc des accusés. On reprocha à la culture française sa
légèreté, son goût pour le brillant superficiel, son obsession de la forme au détriment du fond,
héritages directs d'une tradition pédagogique séculaire dominée par le modèle jésuite et la
rhétorique.\cite{I-1-1-ref1}

Ce rapport se propose d'explorer en profondeur ce moment de basculement, spécifiquement dans le
cadre de l'enseignement du grec ancien. Nous analyserons comment le traumatisme de Sedan a catalysé
le passage d'un \textbf{Modèle Rhétorique}, fondé sur la production (vers, discours) et l'imitation
des modèles (l'\textit{imitatio} jésuite), à un \textbf{Modèle Philologique}, inspiré de
l'Allemagne, fondé sur l'analyse, l'histoire et la grammaire comparée
(\textit{Altertumswissenschaft}). Cette transition, orchestrée par des figures comme Michel Bréal et
Jules Ferry, ne fut pas une simple modernisation, mais une tentative de réarmement moral et
intellectuel de la nation, transformant la classe de grec, jadis lieu de formation du goût et de
l'éloquence, en un laboratoire de la rigueur scientifique.



\subsection{L'ancien régime pédagogique : le modèle rhétorique (jésuite/production)}

Pour saisir la violence de la rupture post-1870, il convient d'abord de dresser l'archéologie du
système qui prévalait jusqu'alors. L'enseignement secondaire français du XIXe siècle restait, dans
ses structures profondes, l'héritier direct des collèges de l'Ancien Régime, et plus spécifiquement
de la pédagogie mise en place par la Compagnie de Jésus à la fin du XVIe siècle.\cite{I-1-1-ref4} Ce
modèle se caractérisait par une finalité précise et une méthode éprouvée : la formation de l'orateur
chrétien par la pratique active de l'éloquence.



\subsection{Le \textit{ratio studiorum} et la primauté de l'\textit{eloquentia}}

La charte fondatrice de cet enseignement était le \textit{Ratio Studiorum} de 1599. Ce document ne
concevait pas l'étude des langues anciennes comme une fin scientifique en soi, mais comme un moyen
de formation morale et esthétique. L'objectif n'était pas de former des érudits capables de
disséquer la langue d'Homère ou de Cicéron, mais de façonner des esprits capables de
\textit{produire} un discours élégant et persuasif. Le latin, et dans une moindre mesure le grec,
étaient des langues vivantes de culture, des outils de communication au sein de la République des
Lettres.\cite{I-1-1-ref4}

La structure même du cursus, menant des classes de Grammaire à celles d'Humanités puis de
Rhétorique, reflétait cette téléologie. L'élève progressait de l'acquisition des règles à
l'imitation des modèles, pour aboutir à la composition originale. Dans ce système, l'exercice roi
n'était pas la traduction vers la langue maternelle (la version), mais la production dans la langue
cible : le \textbf{thème}, le \textbf{discours}, et les \textbf{vers}.\cite{I-1-1-ref6}

La place du grec dans ce dispositif était singulière. Bien que le \textit{Ratio} prescrivît l'étude
du grec, celle-ci demeurait souvent ancillaire par rapport au latin. Le grec était vénéré comme la
source de l'éloquence latine, le modèle indépassable de la beauté formelle, mais il était rarement
étudié pour lui-même avec la rigueur philologique qui caractérisera l'époque suivante. On apprenait
le grec pour mieux comprendre Cicéron, qui avait lui-même imité Démosthène.\cite{I-1-1-ref8} Les
méthodes d'enseignement du grec calquaient celles du latin, privilégiant une approche normative et
esthétique.



\subsection{La logique de la production : vers grecs et discours latins}

Le trait distinctif du modèle rhétorique résidait dans son obsession pour la \og fabrication\fg{}.
L'élève devait faire preuve de sa maîtrise en composant.

\begin{itemize}\item \textbf{Le Discours Latin et Grec :} L'exercice suprême de la classe de Rhétorique consistait à rédiger une harangue en latin (la \textit{concio}), en se glissant dans la peau d'un personnage historique (Hannibal s'adressant à ses troupes, Alexandre le Grand aux mutins). Bien que le discours grec fût plus rare et réservé à l'élite des concours généraux, il participait de la même logique : l'\textit{éthopée}, ou l'art d'imiter les mœurs et le caractère par le style.\cite{I-1-1-ref9}\item \textbf{Les Vers Grecs et Latins :} La composition de poésie, en hexamètres ou en pentamètres, constituait le sommet du raffinement. Il ne s'agissait pas d'être poète par inspiration, mais par technique. L'élève apprenait à assembler des hémistiches, à utiliser des épithètes homériques, à construire des centons (patchworks de vers existants). C'était un exercice de pure virtuosité formelle, très prisé dans les collèges jésuites et conservé par l'Université napoléonienne.\cite{I-1-1-ref9}\item \textbf{Le Thème :} La traduction du français vers le grec (thème grec) était omniprésente. Elle servait de vérification grammaticale mais visait aussi, aux niveaux supérieurs, une qualité stylistique. C'était la clé de voûte des examens, bien plus que la version.\cite{I-1-1-ref11}\end{itemize}Cette pédagogie de la production reposait sur le principe de l'\textbf{Identification}. L'élève ne
devait pas observer l'Antiquité de l'extérieur, comme un objet mort ; il devait l'habiter, la
revivre, parler sa langue. C'était une culture de l'imprégnation et du pastiche.



\subsection{La distinction sociale et l'"honnête homme"}

Sociologiquement, ce modèle remplissait une fonction de distinction bourgeoise essentielle. Comme
l'a analysé Françoise Waquet dans \textit{Le latin ou l'empire d'un signe}, la maîtrise de ces
exercices \og inutiles\fg{} (au sens utilitaire ou économique) fonctionnait comme un marqueur de
classe.\cite{I-1-1-ref14} Savoir tourner un vers latin ou reconnaître une citation de Sophocles
distinguait l'homme bien né, l'\og honnête homme\fg{}, du parvenu ou du technicien.

La \og culture générale\fg{} était définie par cette familiarité avec les classiques, une
familiarité qui se voulait plus esthétique qu'érudite. On cherchait à former le goût, le jugement,
l'élégance morale, plutôt que le savoir positif. Cependant, à mesure que le XIXe siècle avançait
vers l'industrialisation et le scientisme, ce modèle commença à apparaître comme un anachronisme
décoratif, une \og culture de serre\fg{} déconnectée des réalités du monde moderne. C'est sur ce
terreau fragilisé que s'abattit la foudre de 1870.



\subsection{Le traumatisme de sedan et la crise de la conscience française}

La défaite de 1870 ne fut pas perçue comme un simple accident militaire, mais comme le verdict de
l'Histoire jugeant une civilisation. La rapidité de l'effondrement impérial, la capture de Napoléon
III, le siège de Paris, puis la perte des provinces de l'Est plongèrent l'intellectuel français dans
une introspection morbide. Pourquoi la France, patrie des Lumières et des armes, avait-elle été si
aisément surclassée?



\subsection{La "crise allemande" selon claude digeon}

L'analyse de Claude Digeon est ici incontournable. Il décrit comment, après 1870, l'Allemagne devint
pour la France une obsession, un modèle à la fois haï et admiré, une énigme à
déchiffrer.\cite{I-1-1-ref1} Les intellectuels comme Renan, Taine, et plus tard les réformateurs de
l'Instruction publique, diagnostiquèrent que la force de l'Allemagne résidait dans sa méthode, sa
discipline intellectuelle, sa \textit{Wissenschaft}.

L'Allemagne avait su, grâce à la réforme de ses universités par Humboldt et à la rigueur de ses
\textit{Gymnasiums}, créer une élite intellectuelle et militaire rompue à l'analyse, à la précision
et à l'effort. En comparaison, le système français apparut comme une fabrique de \og beaux
parleurs\fg{}, d'esprits brillants mais superficiels, formés à la rhétorique creuse plutôt qu'à la
vérité scientifique. La \og légèreté gauloise\fg{}, jadis vantée, devenait un vice mortel face à la
\og lourdeur\fg{} mais aussi à la \og profondeur\fg{} germanique.\cite{I-1-1-ref1}



\subsection{La critique de la frivolité latine}

Dans ce contexte, les exercices du modèle rhétorique furent désignés comme les boucs émissaires
pédagogiques de la défaite. Le discours latin et les vers grecs devinrent les symboles de cette \og
frivolité\fg{} coupable. À quoi servait-il de savoir imiter le style de Virgile si l'on ignorait la
géographie, les langues vivantes, et la méthode historique critique?

Raoul Frary, dans son pamphlet incendiaire \textit{La question du latin} (1885), cristallisa ces
attaques. Il dénonça le \og fétichisme\fg{} des langues mortes et l'hypocrisie d'un enseignement qui
prétendait former l'esprit par des exercices artificiels dont la plupart des élèves ne tiraient
rien.\cite{I-1-1-ref16} Pour Frary et les \og modernistes\fg{}, la France mourait de son attachement
à une culture de parade. Il fallait remplacer l'élégance par l'efficacité, la rhétorique par la
science.

Le \og Traumatisme de Sedan\fg{} agit ainsi comme un accélérateur de particules historiques. Il
rendit politiquement insoutenable le maintien du statu quo pédagogique. Pour relever la France, il
fallait \og germaniser\fg{} ses méthodes pour mieux combattre l'Allemagne. Il fallait passer du \og
Verbe\fg{} au \og Fait\fg{}.



\subsection{L'importation du modèle philologique : \textit{altertumswissenschaft} et rigueur}

Si le modèle rhétorique était discrédité, par quoi le remplacer? La réponse se trouvait outre-Rhin :
le modèle philologique allemand, ou la \og Science de l'Antiquité\fg{}
(\textit{Altertumswissenschaft}). Ce paradigme, développé dès la fin du XVIIIe siècle par des
figures comme F.A. Wolf, August Boeckh et les frères Humboldt, proposait une approche radicalement
différente des textes anciens.\cite{I-1-1-ref19}



\subsection{La \textit{wissenschaft} contre les belles-lettres}

Contrairement à la tradition française des \og Belles-Lettres\fg{}, qui abordait les textes sous un
angle esthétique et moral, l'\textit{Altertumswissenschaft} allemande considérait l'Antiquité comme
un objet d'étude scientifique total. Le texte n'était plus un modèle à imiter, mais un document à
analyser. Il devait être disséqué à l'aide de l'histoire, de l'archéologie, de l'épigraphie, et
surtout de la linguistique.\cite{I-1-1-ref21}

Dans cette optique, comprendre Homère ne signifiait pas être capable d'écrire des vers à la manière
d'Homère, mais être capable de situer les épopées dans leur contexte historique, d'analyser les
strates dialectales de la langue homérique, et de comprendre la société grecque archaïque. C'était
le passage de l'Antiquité \textit{vécue} (ou rêvée) à l'Antiquité \textit{connue}.\cite{I-1-1-ref20}



\subsection{Le rôle de la grammaire comparée}

Un élément clé de ce nouveau modèle était la linguistique historique et la grammaire comparée. Franz
Bopp avait démontré que le grec et le latin n'étaient pas des isolats sacrés, mais des branches de
la grande famille indo-européenne, régies par des lois phonétiques rigoureuses.\cite{I-1-1-ref22}

Introduire cette perspective dans l'enseignement changeait tout. Le grec n'était plus seulement la
langue des dieux et des muses ; c'était un système linguistique évolutif, comparable au sanskrit ou
au gothique. Cette approche \og scientifique\fg{} de la langue, aride et technique, était perçue par
les réformateurs comme l'antidote idéal à la \og mollesse\fg{} rhétorique française. Elle exigeait
de la précision, de la méthode, et une forme d'ascèse intellectuelle qui rappelait les vertus
prussiennes.\cite{I-1-1-ref22}



\subsection{Les agents du basculement : michel bréal et l'université républicaine}

Ce transfert culturel ne se fit pas tout seul. Il fut l'œuvre d'une génération d'universitaires
français souvent formés en Allemagne, ou du moins fascinés par elle, et qui accédèrent aux postes de
commande sous la IIIe République.



\subsection{Michel bréal : le passeur de la philologie}

La figure tutélaire de ce mouvement est incontestablement Michel Bréal (1832-1915). Professeur au
Collège de France, traducteur de la \textit{Grammaire comparée} de Bopp, Bréal était l'incarnation
de cette synthèse franco-allemande.\cite{I-1-1-ref24}

Dans ses écrits et ses conférences, notamment \textit{De l'enseignement des langues anciennes}
(1891), Bréal mena une guerre sans merci contre la routine pédagogique des lycées.\cite{I-1-1-ref27}
Il fustigea les exercices de production (vers, discours) comme des \og exercices de patience\fg{}
inutiles, qui consommaient un temps précieux sans donner aux élèves une véritable connaissance de
l'Antiquité.

Bréal plaidait pour une pédagogie de la lecture cursive et de l'intelligence. « Il faut lire, lire
beaucoup », répétait-il. L'objectif de l'enseignement du grec devait être de donner un accès direct
à la pensée et à la civilisation grecques, et non de dresser les élèves à des tours de force
grammaticaux. Il voulait remplacer la mémoire par l'intelligence, la répétition par la
compréhension.\cite{I-1-1-ref28}



\subsection{L'institutionnalisation de la recherche}

Sous l'impulsion de Bréal, de Renan, et du ministre Victor Duruy (dès la fin de l'Empire, mais avec
une accélération sous la République), l'enseignement supérieur français se réforma sur le modèle du
séminaire allemand. L'École Pratique des Hautes Études (créée en 1868) fut le fer de lance de cette
nouvelle méthode, privilégiant la recherche par petits groupes, le travail sur manuscrits, et
l'analyse critique, loin des grands cours magistraux d'éloquence de la Sorbonne
d'antan.\cite{I-1-1-ref31}

Les professeurs de lycée, formés dans ces nouvelles moules (notamment via l'Agrégation qui se \og
scientifisa\fg{}), devinrent les vecteurs de cette transformation auprès de la jeunesse. Le
professeur de grec ne devait plus être un rhéteur, mais un savant, un historien de la langue et de
la culture.\cite{I-1-1-ref3}



\subsection{Les champs de bataille de la réforme (1880-1902)}

La transition du modèle Rhétorique au modèle Philologique se concrétisa par une série de réformes
législatives et de débats publics intenses, où se joua le sort des humanités classiques.



\subsection{Le coup de force de 1880 : la mort du discours et du vers}

L'acte fondateur de la rupture fut l'Arrêté du 2 août 1880, porté par Jules Ferry. Ce texte,
véritable déclaration de guerre à la tradition jésuite, supprima purement et simplement les épreuves
de \textbf{discours latin} et de \textbf{vers latins} des examens du baccalauréat et des
concours.\cite{I-1-1-ref9}

C'était une révolution. Du jour au lendemain, l'exercice qui constituait le couronnement des études
depuis trois siècles était aboli. Ferry et ses conseillers (dont Bréal) justifièrent cette mesure
par la nécessité d'\og alléger\fg{} les programmes et de les rendre plus \og sérieux\fg{}.

\begin{itemize}\item La suppression du discours et des vers signifiait la fin de l'idéal de \textbf{Production}. L'État républicain ne demandait plus à ses élites de savoir \og faire du latin\fg{} ou \og faire du grec\fg{}, mais de savoir les \textit{lire} et les \textit{traduire}.\cite{I-1-1-ref10}\item À la place, on institua l'\textbf{explication de texte} et la \textbf{composition française}. L'accent se déplaçait vers la capacité d'analyse critique (en français) des textes anciens. On ne demandait plus à l'élève d'imiter Tite-Live, mais d'expliquer pourquoi Tite-Live écrivait ainsi.\cite{I-1-1-ref9}\end{itemize}

\subsection{La résistance et l'argument de la "gymnastique mentale"}

Face à cette offensive \og utilitariste\fg{} et \og scientiste\fg{}, les défenseurs de la tradition
(souvent catholiques ou conservateurs, mais pas uniquement) se replièrent sur une ligne de défense
désespérée : la théorie de la \textbf{Gymnastique Mentale}.\cite{I-1-1-ref23}

Puisque l'on ne pouvait plus justifier le grec et le latin par leur usage social (on ne parlait plus
latin au Parlement) ni par la production littéraire (désormais interdite), on argua que leur étude
constituait une discipline formelle incomparable pour l'esprit. La difficulté même des déclinaisons
grecques, la rigueur de la syntaxe, \og musclaient\fg{} le cerveau des jeunes bourgeois, leur
donnant une logique et une agilité intellectuelle transférables à toutes les autres
disciplines.\cite{I-1-1-ref23}

C'était un aveu de faiblesse : le contenu culturel du grec importait moins que l'effort fourni pour
l'apprendre. Cet argument, moqué par Raoul Frary comme une \og consolation\fg{} pour un enseignement
inefficace 18, devint pourtant le credo officiel de l'enseignement secondaire classique jusqu'au
milieu du XXe siècle.



\subsection{La réforme de 1902 : la consécration de la bifurcation}

L'aboutissement de ce processus fut la grande réforme de 1902. Elle créa des filières distinctes
dans le secondaire, mettant fin au monopole des humanités classiques. Désormais, il existait un
baccalauréat \og Moderne\fg{} (Sciences et Langues vivantes) à égalité de dignité avec le
baccalauréat \og Classique\fg{} (Latin-Grec).\cite{I-1-1-ref37}

Pour le grec, ce fut le coup de grâce de son statut universel. Il devint une option de distinction
pour l'élite de la section A. Mais surtout, la réforme entérina la victoire de la méthode
philologique (dite \og méthode de traduction\fg{}) pour les langues anciennes, par opposition à la
\og méthode directe\fg{} introduite pour les langues vivantes.\cite{I-1-1-ref38} Les langues
vivantes servaient à communiquer (parler, écouter) ; les langues anciennes servaient à analyser
(traduire, commenter). La dichotomie était scellée.



\subsection{La transformation de l'enseignement du grec : de l'art à la science}

Concrètement, qu'est-ce que cela changea dans la classe de grec? Le passage du modèle Rhétorique au
modèle Philologique modifia la nature même des exercices et la posture de l'élève.



\subsection{La survie du thème et l'hégémonie de la version}

Si le discours et les vers disparurent (sauf comme curiosités), le \textbf{thème grec} (traduction
du français vers le grec) survécut, notamment dans les classes supérieures et les concours
d'excellence (Agrégation, Khâgne).\cite{I-1-1-ref13}

Cependant, sa nature changea. Dans le modèle jésuite, le thème était un exercice de style, une étape
vers la création libre. Dans le modèle philologique, il devint un instrument de contrôle
grammatical. Il servait à vérifier que l'élève avait parfaitement assimilé les règles de la syntaxe
et de la morphologie. La \og beauté\fg{} du thème comptait moins que son \og exactitude\fg{}
philologique. Le thème devint une épreuve technique, un \og verrou\fg{} de sécurité pour
l'agrégation, garantissant que le futur professeur possédait la science de la
langue.\cite{I-1-1-ref11}

Parallèlement, la \textbf{version grecque} (traduction du grec vers le français) devint l'exercice
roi de l'intelligence. Elle exigeait de comprendre le texte dans ses nuances, de l'analyser
logiquement, et de le rendre dans un français précis. C'était l'exercice d'analyse par excellence,
conforme à l'esprit critique républicain.\cite{I-1-1-ref11}



\subsection{L'histoire littéraire et la civilisation}

Le cours de grec s'enrichit de nouvelles matières. On ne lisait plus seulement les textes ; on
étudiait l'histoire de la littérature grecque, les institutions d'Athènes, l'art grec. Les manuels
se remplirent de gravures de monuments et de statues, reflet de l'essor de
l'archéologie.\cite{I-1-1-ref2} L'objectif était de donner une \og culture\fg{} historique, de
comprendre le \og génie grec\fg{} dans son ensemble, tel que le définissaient Taine ou Renan, et non
plus seulement d'apprécier une belle tournure de phrase.

C'était une vision plus riche, plus intellectuelle, mais aussi plus froide. Le texte grec était mis
à distance. Il devenait un objet d'étude, séparé de l'élève par l'épaisseur des siècles et la
rigueur de la science. L'identification émotionnelle et mimétique du modèle rhétorique s'effaçait
devant l'objectivité du savant.

Le basculement de 1870-1902, du modèle Rhétorique au modèle Philologique, fut une réponse directe au
traumatisme de la défaite. Pour relever la France, l'Université républicaine choisit d'adopter les
armes intellectuelles du vainqueur : la science, l'analyse, l'histoire critique. Elle sacrifia la
vieille pédagogie jésuite de l'éloquence et de la production, jugée coupable d'avoir amolli les
esprits, sur l'autel de la rigueur philologique.

Ce tournant permit indéniablement de moderniser l'enseignement français et de former des générations
de savants de haut niveau. La France rattrapa son retard scientifique sur l'Allemagne dans le
domaine des études anciennes. Cependant, on peut se demander si cette victoire ne fut pas payée au
prix fort.

En transformant le grec en une discipline purement analytique et \og gymnastique\fg{}, en le coupant
de la créativité et du plaisir de \og faire\fg{}, les réformateurs ont peut-être contribué à tarir
sa sève. Devenu une \og science de l'Antiquité\fg{} aride, justifiée par la seule formation de
l'esprit, le grec perdit peu à peu son attrait auprès des élèves et des familles, préparant son lent
déclin au XXe siècle. La \og Crise allemande\fg{} avait été surmontée, mais le \og Miracle
grec\fg{}, lui, avait été mis en bouteille dans le formol de la philologie.



\subsection{Annexes et données structurées}



\subsection{Tableau 1 : confrontation des deux modèles pédagogiques}

\begin{table}[h!]\small\centering\begin{tabular}{| p{2.3cm} | p{3.5cm} | p{3.5cm} |}\hline\textbf{Caractéristique} & \textbf{Modèle Rhétorique (Jésuite / Avant 1870)} & \textbf{Modèle Philologique (Allemand / Après 1880)} \\ \hline\textbf{Objectif Ultime} & \textit{Eloquentia} (Éloquence), Formation Morale & \textit{Science} (Savoir Positif), Compréhension Historique \\ \hline\textbf{Exercices Clés} & Discours, Vers, Thème (Production) & Explication de texte, Version, Histoire (Analyse) \\ \hline\textbf{Rapport au Texte} & Imitation / Identification / \og Habiter le texte\fg{} & Analyse / Objectivation / \og Comprendre le texte\fg{} \\ \hline\textbf{Statut du Grec} & Modèle esthétique, serviteur du Latin & Objet d'\textit{Altertumswissenschaft}, donnée linguistique \\ \hline\textbf{Idéal de l'Élève} & L'Orateur Chrétien / L'Honnête Homme & Le Philologue / L'Esprit Critique / Le Citoyen \\ \hline\textbf{Influence Majeure} & \textit{Ratio Studiorum} (Compagnie de Jésus) & Universités Allemandes (Humboldt/Wolf/Boeckh) \\ \hline\end{tabular}\end{table}

\subsection{Tableau 2 : chronologie du basculement}

\begin{table}[h!]\small\centering\begin{tabular}{| p{2.3cm} | p{3.5cm} | p{3.5cm} |}\hline\textbf{Date} & \textbf{Événement} & \textbf{Impact sur l'Enseignement du Grec} \\ \hline\textbf{1870} & \textbf{Défaite de Sedan} & Choc moral. Prise de conscience de la supériorité du modèle éducatif prussien. \\ \hline\textbf{1872} & M. Bréal, \textit{Quelques mots sur l'instruction publique} & Appel à l'adoption des méthodes scientifiques et critiques allemandes. \\ \hline\textbf{1880} & \textbf{Arrêté du 2 août (Jules Ferry)} & \textbf{Suppression officielle du Discours Latin et des Vers.} Le tournant décisif. \\ \hline\textbf{1885} & Raoul Frary, \textit{La question du latin} & Attaque publique contre l'utilité des exercices classiques (\og Fétichisme\fg{}). \\ \hline\textbf{1890} & Instructions de 1890 & Consécration de l'argument de la \og Gymnastique Mentale\fg{} pour sauver les humanités. \\ \hline\textbf{1891} & Bréal, \textit{De l'enseignement des langues anciennes} & Manifeste pour la méthode philologique (lecture vs par cœur). \\ \hline\textbf{1902} & \textbf{Réforme de 1902} & Création des filières modernes. Le grec devient une spécialité optionnelle d'élite. \\ \hline\end{tabular}\end{table}

\subsection{Tableau 3 : l'évolution des épreuves (exemple de l'agrégation et du baccalauréat)}

\begin{table}[h!]\small\centering\begin{tabular}{| p{2.3cm} | p{3.5cm} | p{3.5cm} |}\hline\textbf{Période} & \textbf{Épreuves Écrites Dominantes} & \textbf{Nature de l'Oral} \\ \hline\textbf{Pré-1880} & Discours Latin, Vers Latins/Grecs, Thème Grec & Explication rhétorique (appréciation du style, des figures) \\ \hline\textbf{Post-1880} & Version (Grec/Latin), Thème (Grammatical), Dissertation Française & Explication philologique (grammaire, histoire, contexte) \\ \hline\textbf{Analyse} & Disparition de la créativité au profit de la vérification technique et de la capacité de traduction. & L'oral devient un cours d'histoire et de grammaire appliquée. \\ \hline\end{tabular}\end{table}